% ============================================================
% Section 33: 一维离子晶格的马德隆常数
% ============================================================
\section{固体理论：一维离子晶格的马德隆常数}
\label{sec:madelung-1d}

\begin{modelbox}
设一晶体由正负一价离子交错等距排列组成一维晶格。证明其马德隆（Madelung）常数为
\[
  M = 2\ln 2.
\]
\end{modelbox}

\vspace{0.3em}
\begin{insightbox}{结论与数量级}
\[
  \boxed{M = 2\ln 2 \approx 1.386}.
\]
作为对比，三维 NaCl 结构的马德隆常数 $M_{\text{NaCl}} \approx 1.748$。一维链的 $M$ 更小，根本原因是\textbf{一维中每个离子的最近邻异号离子太少}。

\textbf{直观对比：}
\begin{center}
\renewcommand{\arraystretch}{1.4}
\begin{tabular}{|l|c|c|}
\hline
\textbf{维度} & \textbf{最近邻异号离子数} & \textbf{马德隆常数} \\
\hline
一维链 & 2（左、右各1个） & $1.386$ \\
三维 NaCl & 6（前、后、左、右、上、下） & $1.748$ \\
\hline
\end{tabular}
\end{center}

三维中你被\textbf{6个异号离子}近距离包围，吸引力更强；一维中只有\textbf{2个}。虽然远处的同号离子会排斥你，但它们距离更远（按 $r_0$ 的倍数），且被异号离子``挡在中间''，排斥效果被部分抵消。最近邻的数量优势是决定 $M$ 大小的核心因素。
\end{insightbox}

\subsection{固体物理题型反射弧——考场不慌指南}

\begin{modelbox}{看到什么关键词$\to$启动什么公式}
\begin{center}
\renewcommand{\arraystretch}{1.6}
\begin{tabular}{|p{4.2cm}|p{4.5cm}|p{5.5cm}|}
\hline
\textbf{题目关键词} & \textbf{题型判断} & \textbf{核心操作} \\
\hline
``相邻原子间相互作用势'' ``平衡距离'' ``晶格能'' & 晶体结合能（LJ/幂律势） & 一阶导$=0$求平衡，代入求能，二阶导乘键长得弹性模量 \\
\hline
``一维离子晶格'' ``马德隆常数'' ``库仑能'' & 马德隆常数 & 写交错级数$\sum (\pm 1)/n$，利用$\ln(1+x)$或$\ln 2$求和 \\
\hline
``一维单原子链'' ``声子'' ``色散关系'' & 晶格振动 & 简谐近似$+$牛顿方程$+$Bloch行波解$\to$求$\omega(q)$ \\
\hline
``玻尔兹曼方程'' ``弛豫时间'' ``电场中分布函数'' & 固体输运 & 稳态方程$+$按电场幂次展开$+$识别平移算符 \\
\hline
``约瑟夫森结'' ``SQUID'' ``磁通量子'' & 超导器件 & 约瑟夫森方程$I=I_c\sin\varphi$ $+$ 相位积累$+$磁通量子化 \\
\hline
``能带'' ``周期性势场'' ``Bloch波'' & 能带论 & 近自由电子$/$紧束缚模型，求色散$E(k)$，算有效质量 \\
\hline
\end{tabular}
\end{center}
\end{modelbox}

\begin{tipbox}{固体物理的``三板斧''}
无论遇到哪种题型，底层逻辑几乎都可以归为以下三步：
\begin{enumerate}[nosep]
  \item \textbf{写势能/哈密顿量}：明确相互作用形式（LJ、库仑、谐振子、紧束缚等）；
  \item \textbf{求平衡/本征态}：导数为零、对角化、解运动方程；
  \item \textbf{求响应/修正}：二阶导（弹性模量）、微扰（输运）、色散关系（声子/能带）。
\end{enumerate}
固体物理考题的吓人之处往往在于符号多、近似多，但核心结构高度重复。只要认准题型，套对应的三板斧，大部分分数都能拿下。
\end{tipbox}

\subsection{概念铺垫：什么是马德隆常数？}

\begin{modelbox}{库仑相互作用能——两个点电荷的``共享能量''}
两个点电荷 $q_1$, $q_2$ 相距 $r$，它们的库仑相互作用能为
\[
  u(r) = \frac{q_1 q_2}{4\pi\varepsilon_0 r}.
\]
\begin{itemize}[nosep]
  \item \textbf{异号电荷}（$q_1 q_2 < 0$）：$u < 0$，能量为\textbf{负}，表示\textbf{吸引}——系统能量降低，更稳定；
  \item \textbf{同号电荷}（$q_1 q_2 > 0$）：$u > 0$，能量为\textbf{正}，表示\textbf{排斥}——系统能量升高，不稳定。
\end{itemize}
在离子晶体中，每个离子同时与所有其他离子有库仑作用。\textbf{总库仑能}是所有不同离子对的 $u(r)$ 之和。
\end{modelbox}

\begin{modelbox}{计数方式：为什么总能量不是``人数$\times$个人能量''？}
晶体中有 $N$ 对离子（$N$ 个正离子 + $N$ 个负离子，共 $2N$ 个）。选取一个正离子为参考，它与其他所有离子的相互作用能之和记为 $u_0$。

\textbf{陷阱}：若直接把 $u_0$ 乘以总离子数 $2N$，得到 $(2N)u_0$——这\textbf{错了}，因为每对相互作用被算了\textbf{两次}（$i$ 算 $j$ 时一次，$j$ 算 $i$ 时一次）。

\textbf{正确做法}：
\[
  U = \underbrace{\frac{1}{2}}_{\text{除重复}} \times \underbrace{(2N)}_{\text{总离子数}} \times u_0
    = N u_0.
\]
因此标准形式写为
\[
  U = -\frac{N M e^2}{4\pi\varepsilon_0 r_0},
\]
其中 $N$ 是\textbf{离子对数}（也是正离子数或负离子数），$M$ 是无量纲的马德隆常数。
\end{modelbox}

\begin{insightbox}{为什么不直接求和？——长程灾难}
库仑势 $1/r$ 是\textbf{长程}的，衰减极慢。在一维链中，第 $n$ 个邻居的贡献只按 $1/n$ 衰减。若所有离子都是同号，级数 $\sum 1/n$ 会发散；幸运的是，交替电荷使级数变成\textbf{交错级数}，正负抵消后收敛到有限值。
\end{insightbox}

\subsection{马德隆常数的推导}

\subsubsection{选取参考离子并写级数}

选取晶格中某一正离子为\textbf{参考离子}（位于原点），设最近邻间距为 $r_0$。第 $n$ 个离子位于 $nr_0$，电荷为 $(-1)^n e$（$n=1$ 为负离子，$n=2$ 为正离子，依此类推）。

参考离子与右侧第 $n$ 个离子的库仑能为
\[
  u_n = \frac{e \cdot (-1)^n e}{4\pi\varepsilon_0 \cdot nr_0}
      = \frac{(-1)^n e^2}{4\pi\varepsilon_0 nr_0}.
\]

右侧所有离子对参考离子的总能量：
\[
  u_{\text{右}} = \sum_{n=1}^{\infty} u_n
               = \frac{e^2}{4\pi\varepsilon_0 r_0} \sum_{n=1}^{\infty} \frac{(-1)^n}{n}.
\]

由于晶格左右对称，左侧贡献完全相同：
\[
  u_{\text{左}} = u_{\text{右}}.
\]

因此，参考离子的总库仑相互作用能：
\[
  u_0 = u_{\text{左}} + u_{\text{右}}
      = \frac{2e^2}{4\pi\varepsilon_0 r_0} \sum_{n=1}^{\infty} \frac{(-1)^n}{n}. \tag{1}
\]

\subsubsection{级数求和}

利用自然对数的泰勒展开：
\[
  \ln(1+x) = x - \frac{x^2}{2} + \frac{x^3}{3} - \frac{x^4}{4} + \cdots
            = \sum_{n=1}^{\infty} (-1)^{n-1} \frac{x^n}{n},
  \quad (-1 < x \le 1).
\]

令 $x=1$（收敛域边界，交错级数收敛）：
\[
  \ln 2 = 1 - \frac{1}{2} + \frac{1}{3} - \frac{1}{4} + \cdots
        = \sum_{n=1}^{\infty} \frac{(-1)^{n-1}}{n}.
\]

注意式 (1) 中级数的符号：
\[
  \sum_{n=1}^{\infty} \frac{(-1)^n}{n}
  = -1 + \frac{1}{2} - \frac{1}{3} + \frac{1}{4} - \cdots
  = -\sum_{n=1}^{\infty} \frac{(-1)^{n-1}}{n}
  = -\ln 2.
\]

代入 (1)：
\[
  u_0 = \frac{2e^2}{4\pi\varepsilon_0 r_0} (-\ln 2)
      = -\frac{(2\ln 2)\, e^2}{4\pi\varepsilon_0 r_0}. \tag{2}
\]

\subsubsection{提取马德隆常数}

晶体中有 $N$ 对离子（$2N$ 个离子）。总库仑能为
\[
  U = \underbrace{\frac{1}{2}(2N)}_{=N} u_0 = -\frac{N (2\ln 2)\, e^2}{4\pi\varepsilon_0 r_0}.
\]

对比标准形式 $U = -\dfrac{N M e^2}{4\pi\varepsilon_0 r_0}$（$N$ 为离子对数），直接读出
\[
  \boxed{ M = 2\ln 2 }.
\]

\begin{tipbox}{符号约定的自洽性检验}
式 (2) 中 $u_0 < 0$，表示库仑作用总体为吸引——这与离子晶体的物理图像一致（正负离子交替，异号占优）。若符号算反了，会得到 $M < 0$，意味着晶体因库仑排斥而解体，显然荒谬。
\end{tipbox}

\subsection{物理图像与拓展}

\subsubsection{一维 vs 三维：维度对 $M$ 的影响}

\begin{center}
\renewcommand{\arraystretch}{1.5}
\begin{tabular}{|l|c|l|}
\hline
\textbf{维度} & \textbf{马德隆常数} & \textbf{级数衰减} \\
\hline
一维交替链 & $M = 2\ln 2 \approx 1.386$ & $1/n$（慢，但交错收敛） \\
二维 NaCl 型 & $M \approx 1.615$ & $1/r$（更慢，需 Ewald 加速） \\
三维 NaCl & $M \approx 1.748$ & $1/r$（直接求和发散，必须用 Ewald） \\
三维 CsCl & $M \approx 1.763$ & $1/r$（同上） \\
\hline
\end{tabular}
\end{center}

维度越高，$M$ 越大。原因是三维中每个离子被更多异号离子包围，同号离子的排斥被更有效地屏蔽。

\subsubsection{三维中的埃瓦尔德（Ewald）求和——为什么需要它？}

\begin{modelbox}{三维求和的困境}
三维中直接计算 $\sum (\pm 1)/r$ 会遇到两个致命问题：
\begin{enumerate}[nosep]
  \item \textbf{条件收敛}：级数的和依赖于求和顺序（按球壳、立方壳等不同截断方式给出不同结果）；
  \item \textbf{表面效应}：球壳截断时，表面有净电荷（未完全配对），产生偶极矩，导致结果依赖于截断半径。
\end{enumerate}
一维链没有表面效应问题，因为``左右各一条路''，正负严格配对。
\end{modelbox}

\begin{insightbox}{Ewald 的核心思想：给每个离子穿一件``羽绒服''}
想象每个点电荷不是无限尖锐的点，而是被一团``高斯分布''的蓬松电荷云包围（像羽绒服）。

\textbf{拆分策略}：把 $1/r$ 拆成两部分
\[
  \frac{1}{r} = \underbrace{\frac{\text{erfc}(\eta r)}{r}}_{\text{实空间：短程尖锐}} + \underbrace{\frac{\text{erf}(\eta r)}{r}}_{\text{倒空间：长程平滑}}.
\]

其中 $\text{erf}(x)$ 是误差函数（$\text{erf}(0)=0$, $\text{erf}(\infty)=1$），$\text{erfc}(x)=1-\text{erf}(x)$ 是互补误差函数。

\textbf{第一项（实空间部分）}：$\text{erfc}(\eta r)/r$
\begin{itemize}[nosep]
  \item 当 $r$ 很小时，$\text{erfc}(\eta r)\approx 1$，这一项接近 $1/r$，还原了近距离的尖锐作用；
  \item 当 $r$ 很大时，$\text{erfc}(\eta r)$ 按 $e^{-\eta^2 r^2}$ 指数衰减到 0，\textbf{非常快}；
  \item 因此实空间中只需计算最近邻少数几个离子的贡献。
\end{itemize}

\textbf{第二项（倒空间部分）}：$\text{erf}(\eta r)/r$
\begin{itemize}[nosep]
  \item 在实空间中，这一项变化缓慢（长程），但做傅里叶变换（转到倒空间）后，只包含少量低波矢成分；
  \item 在倒空间中，只需要取前几项倒格矢 $\mathbf{G}$ 就能精确描述；
  \item 因此倒空间中也只需计算少数几项。
\end{itemize}

\textbf{可调参数 $\eta$ 的作用}：
\begin{itemize}[nosep]
  \item $\eta$ 大：高斯峰尖锐，实空间收敛极快，但倒空间收敛慢；
  \item $\eta$ 小：高斯峰平坦，实空间收敛慢，但倒空间收敛快；
  \item 最优策略：取适中的 $\eta$，使两部分计算量大致相等。
\end{itemize}
\end{insightbox}
